
\section{Planteamiento científico}

La pregunta tradicional sobre una estructura histórica suele formularse en sentido directo: qué mecanismo podría haber producido el objeto observado. Ese enfoque es legítimo, pero muchas historias diferentes pueden ser físicamente plausibles y producir estados terminales semejantes. La plausibilidad no equivale a identificación histórica.

Sea \(X_T\) el estado terminal \emph{as-built}, \(Y_{\mathrm{now}}\) el conjunto de observaciones modernas y \(\mathcal L\) la familia de leyes físicas y modelos admisibles. El objetivo de IRL es inferir

\begin{equation}
\boxed{
P(H,\Theta\mid Y_{\mathrm{now}},\mathcal L)
}
\end{equation}

donde \(H\) representa historias constructivas discretas o parcialmente ordenadas y \(\Theta\) reúne variables continuas como fechas, densidades, rutas, fuerzas, tolerancias, parámetros materiales y geometrías latentes.

\begin{principle}[No unicidad por defecto]
Si dos o más historias permanecen observacionalmente indistinguibles bajo los operadores de observación y la evidencia disponible, el resultado correcto es una clase de equivalencia o una distribución posterior, no una narración única.
\end{principle}

La Gran Pirámide de Khufu constituye el laboratorio arqueológico principal por combinar geometría monumental, diferencias litológicas, cámaras y corredores, canteras, lecho rocoso, problemas de transporte, estructura interna, observaciones muográficas y una cronología constructiva no directamente observada \cite{petrie1883,arnold1991,lehner1997,morishima2017,procureur2023,aeraquarry,ghoneim2024}.

\section{Dos problemas inversos encadenados}

La pirámide observable hoy no es idéntica al objeto que existía al finalizar su construcción. Pérdida de revestimiento, erosión, fracturas, extracción de piedra, excavación y reparaciones alteraron el sistema.

Definimos

\begin{equation}
X_{\mathrm{now}}
=
\post(X_T;\Theta_{\mathrm{post}})
\end{equation}

y

\begin{equation}
Y_{\mathrm{now}}
=
\obs[X_{\mathrm{now}}]+\varepsilon.
\end{equation}

El estado terminal construido surge de

\begin{equation}
X_T
=
\build(C,D,\Gamma;\Theta_{\mathrm{build}}),
\end{equation}

donde \(C\) representa restricciones de concepción, \(D\) información de diseño y \(\Gamma\) la historia constructiva.

Por tanto,

\begin{align}
\mathcal I_{\mathrm{post}} &: 
Y_{\mathrm{now}}
\rightarrow
P(X_T\mid Y_{\mathrm{now}}),\\
\mathcal I_{\mathrm{build}} &: 
X_T
\rightarrow
P(\Gamma,D,C\mid X_T,E,\mathcal L).
\end{align}

Esta separación evita atribuir al proceso original propiedades que pertenecen a la historia postconstructiva.

\section{Tiempo inverso: coordenada inferencial}

Sea \(t\in[0,T]\) el tiempo constructivo directo. Definimos

\begin{equation}
\boxed{\tau=T-t.}
\end{equation}

Entonces

\begin{equation}
\frac{d}{d\tau}
=
-\frac{d}{dt},
\qquad
\frac{d^2}{d\tau^2}
=
\frac{d^2}{dt^2}.
\end{equation}

La segunda igualdad no implica reversibilidad termodinámica. Fricción, daño y disipación destruyen información microscópica. IRL reconstruye antecedentes que, evolucionados hacia adelante bajo física ordinaria, son compatibles con el estado observado.

\begin{remark}
El desmontaje inverso es un algoritmo de inferencia y un orden de estados candidatos. No afirma que un proceso histórico pueda reproducirse físicamente al revés.
\end{remark}

\section{Geometría ideal y relojes no equivalentes}

Consideremos una pirámide cuadrada ideal de lado de base \(a\), altura \(H\), coordenada vertical \(z\) y altura normalizada \(u=z/H\).

El área de una sección horizontal es

\begin{equation}
A(z)
=
a^2
\left(1-\frac{z}{H}\right)^2.
\end{equation}

\begin{proposition}[Volumen ideal]
\begin{equation}
\boxed{
V=\frac{a^2H}{3}.
}
\end{equation}
\end{proposition}

Bajo densidad uniforme, la fracción de masa situada por encima de una cota \(u\) es

\begin{equation}
\boxed{
\Gamma_M(u)
=
(1-u)^3.
}
\end{equation}

Así,

\begin{equation}
\Gamma_M(1/2)=1/8=12.5\%.
\end{equation}

La energía potencial respecto de la base es

\begin{equation}
U
=
\rho g
\int_0^H zA(z)\,dz
=
\frac14MgH.
\end{equation}

La fracción de potencial situada por encima de \(u\) es

\begin{equation}
\boxed{
\Gamma_G(u)
=
1-6u^2+8u^3-3u^4.
}
\end{equation}

En \(u=1/2\),

\begin{equation}
\Gamma_G(1/2)=5/16=31.25\%.
\end{equation}

\begin{principle}[No equivalencia de relojes]
\begin{equation}
\boxed{
\tau
\neq
z/H
\neq
\Gamma_M
\neq
\Gamma_G.
}
\end{equation}
La cronología no puede sustituirse por una sola variable geométrica.
\end{principle}
